% ---------------------------------------------------------------------------
%  Preambolo comune del documento
% ---------------------------------------------------------------------------
\usepackage[T1]{fontenc}
\usepackage[utf8]{inputenc}
\usepackage[italian]{babel}

\usepackage{lmodern}
\usepackage{textcomp}
\usepackage{microtype}

\usepackage[a4paper,top=2.6cm,bottom=2.7cm,inner=2.7cm,outer=2.2cm,
            headsep=14pt]{geometry}

\usepackage{amsmath}
\usepackage{amssymb}
\usepackage{mathtools}
\usepackage{graphicx}
\usepackage{xcolor}
\usepackage{booktabs}
\usepackage{longtable}
\usepackage{array}
\usepackage{enumitem}
\usepackage{fancyhdr}
\usepackage{titlesec}
\usepackage{caption}
\usepackage{listings}
\usepackage{tikz}
\usetikzlibrary{arrows.meta,positioning,calc,shapes.geometric}
\usepackage[hidelinks,pdfusetitle]{hyperref}

% --- tavolozza --------------------------------------------------------------
\definecolor{inchiostro}{HTML}{1B1F23}
\definecolor{blunotte}{HTML}{1F4E79}
\definecolor{rossoterra}{HTML}{B3352B}
\definecolor{verdesalvia}{HTML}{2A7F62}
\definecolor{grigiocarta}{HTML}{F4F5F7}
\definecolor{grigiotratto}{HTML}{C8CCD2}
\definecolor{grigiotesto}{HTML}{6B7280}

\color{inchiostro}

% --- titolazione ------------------------------------------------------------
\titleformat{\chapter}[display]
  {\normalfont\sffamily\huge\bfseries\color{blunotte}}
  {\normalsize\sffamily\mdseries\color{grigiotesto}%
   \MakeUppercase{\chaptertitlename}\ \thechapter}
  {8pt}{\Huge}
\titlespacing*{\chapter}{0pt}{-10pt}{26pt}

\titleformat{\section}
  {\normalfont\sffamily\large\bfseries\color{blunotte}}{\thesection}{0.7em}{}
\titleformat{\subsection}
  {\normalfont\sffamily\normalsize\bfseries\color{inchiostro}}
  {\thesubsection}{0.6em}{}

\setcounter{secnumdepth}{2}
\setcounter{tocdepth}{1}

% --- testatine --------------------------------------------------------------
\pagestyle{fancy}
\fancyhf{}
\renewcommand{\headrulewidth}{0.4pt}
\renewcommand{\headrule}{\hbox to\headwidth{\color{grigiotratto}%
  \leaders\hrule height \headrulewidth\hfill}}
\fancyhead[LE]{\small\sffamily\color{grigiotesto}\leftmark}
\fancyhead[RO]{\small\sffamily\color{grigiotesto}\rightmark}
\fancyfoot[LE,RO]{\small\sffamily\thepage}
\renewcommand{\chaptermark}[1]{\markboth{#1}{}}
\renewcommand{\sectionmark}[1]{\markright{#1}}

\fancypagestyle{plain}{\fancyhf{}%
  \renewcommand{\headrulewidth}{0pt}%
  \fancyfoot[LE,RO]{\small\sffamily\thepage}}

% --- didascalie -------------------------------------------------------------
\captionsetup{font=small,labelfont={sf,bf,color=blunotte},
              labelsep=period,width=0.92\linewidth}

% --- ambiente per le espressioni --------------------------------------------
\newcommand{\verificata}{%
  \textcolor{verdesalvia}{\ensuremath{\checkmark}}}

\newcommand{\esprheader}[3]{%
  \par\addvspace{7pt}\noindent
  {\footnotesize\sffamily
   \textcolor{blunotte}{\textbf{\##1}}\,#3\hfill
   \textcolor{grigiotesto}{#2}}%
  \par\nobreak\vspace{-6pt}}

\newenvironment{espr}[3]
  {\esprheader{#1}{#2}{#3}\abovedisplayskip=6pt\belowdisplayskip=6pt}
  {\par\addvspace{2pt}}

% le catene di uguaglianze possono essere lunghe: si consente la rottura
\allowdisplaybreaks[3]

% --- intestazione di ogni quaderno ------------------------------------------
\newcommand{\quaderno}[3]{%
  \begingroup
  \setlength{\fboxsep}{9pt}%
  \noindent\colorbox{grigiocarta}{%
    \begin{minipage}{\dimexpr\linewidth-2\fboxsep\relax}
      \footnotesize\sffamily\color{grigiotesto}
      \textbf{Sorgente} \texttt{#1}\hfill
      \textbf{Salvato il} #2\hfill
      \textbf{Espressioni} #3
    \end{minipage}}%
  \endgroup
  \par\addvspace{14pt}}

% --- riquadro per le osservazioni -------------------------------------------
\newenvironment{osservazione}[1][Osservazione]
  {\par\addvspace{8pt}\begingroup
   \setlength{\leftskip}{1.2em}%
   \noindent{\sffamily\bfseries\color{rossoterra}#1.}\ \ignorespaces}
  {\par\endgroup\addvspace{8pt}}

% --- codice sorgente --------------------------------------------------------
\lstdefinestyle{pitone}{
  language=Python,
  basicstyle=\ttfamily\scriptsize,
  keywordstyle=\color{blunotte}\bfseries,
  commentstyle=\color{verdesalvia}\itshape,
  stringstyle=\color{rossoterra},
  numbers=left, numberstyle=\tiny\color{grigiotesto}, numbersep=8pt,
  showstringspaces=false,
  breaklines=true, breakatwhitespace=true,
  frame=leftline, framerule=0.6pt, rulecolor=\color{grigiotratto},
  xleftmargin=16pt, framexleftmargin=16pt,
  extendedchars=true, inputencoding=utf8,
  literate={à}{{\`a}}1 {è}{{\`e}}1 {é}{{\'e}}1 {ì}{{\`i}}1 {ò}{{\`o}}1
           {ù}{{\`u}}1 {'}{{\textquotesingle}}1
}
\lstset{style=pitone}

% --- varie ------------------------------------------------------------------
\newcommand{\dfw}[1]{\texttt{#1}}
\newcommand{\der}[1]{\texttt{\small #1}}
\renewcommand{\arraystretch}{1.15}
\setlength{\parskip}{5pt}
\setlength{\parindent}{0pt}
\setlength{\emergencystretch}{2em}
\hbadness=10000
\vbadness=10000
